% ============================================================
%  پروپوزال رویداد مخابرات کوانتومی — همراه اول
%  کامپایلر: XeLaTeX
% ============================================================
\documentclass[a4paper,12pt]{article}

% ── بسته‌های اصلی ──────────────────────────────────────────
\usepackage{fontspec}

\usepackage{geometry}
\usepackage{xcolor}
\usepackage{graphicx}
\usepackage{tikz}
\usepackage{pgffor}
\usepackage{array}
\usepackage{tabularx}
\usepackage{booktabs}
\usepackage{multirow}
\usepackage{enumitem}
\usepackage{mdframed}
\usepackage{titlesec}
\usepackage{fancyhdr}
\usepackage{hyperref}
\usepackage{eso-pic}   % برای پس‌زمینه صفحه

\usetikzlibrary{shapes.geometric, arrows.meta, positioning, decorations.pathreplacing, calc}

% ── تنظیمات صفحه ──────────────────────────────────────────
\geometry{
  a4paper,
  top=2.5cm, bottom=2.5cm,
  left=2.5cm, right=2.5cm,
  headheight=1.5cm
}

\usepackage{xepersian}

% ── فونت‌ها ────────────────────────────────────────────────
% فونت پیش‌فرض: XB Niloofar (اگر نصب نیست، IranNastaliq یا هر فونت فارسی موجود)
% برای تغییر، نام فونت را در خطوط زیر عوض کنید
\settextfont{XB Niloofar}
\setdigitfont{XB Niloofar}

% فونت انگلیسی
\newfontfamily\enfont{TeX Gyre Termes}

% ── پالت رنگ (بر پایه هویت همراه اول: آبی عمیق + جزئیات کوانتوم) ─
\definecolor{mci-blue}{HTML}{003A70}      % آبی اصلی همراه اول
\definecolor{mci-accent}{HTML}{0074C8}    % آبی روشن‌تر
\definecolor{mci-gold}{HTML}{C8973A}      % طلایی برای تأکید
\definecolor{mci-light}{HTML}{E8F1FA}     % پس‌زمینه روشن
\definecolor{mci-dark}{HTML}{1A1A2E}      % تیره برای متن روی رنگ
\definecolor{mci-gray}{HTML}{5A6A7A}      % خاکستری متن ثانویه
\definecolor{mci-border}{HTML}{B0C4D8}    % رنگ خطوط

% ── تنظیمات هایپرلینک ─────────────────────────────────────
\hypersetup{
  colorlinks=true,
  linkcolor=mci-blue,
  urlcolor=mci-accent,
  pdfauthor={تیم تحقیق و فناوری همراه اول},
  pdftitle={پروپوزال رویداد مخابرات کوانتومی}
}

% ── سرصفحه و پاصفحه ──────────────────────────────────────
\pagestyle{fancy}
\fancyhf{}
\renewcommand{\headrulewidth}{0.4pt}
\renewcommand{\headrule}{\hbox to\headwidth{\color{mci-blue}\leaders\hrule height \headrulewidth\hfill}}
\fancyhead[R]{\small\color{mci-gray} تیم تحقیق و فناوری — همراه اول}
\fancyhead[L]{\small\color{mci-gray} پروپوزال رویداد کوانتوم}
\fancyfoot[C]{\small\color{mci-gray}\thepage}

% ── عنوان‌بندی بخش‌ها ─────────────────────────────────────
\titleformat{\section}
  {\large\bfseries\color{mci-blue}}
  {\thesection.}{0.5em}{}
  [\textcolor{mci-accent}{\titlerule[1.5pt]}]

\titleformat{\subsection}
  {\normalsize\bfseries\color{mci-dark}}
  {\thesubsection.}{0.5em}{}

\setlength{\parindent}{0pt}
\setlength{\parskip}{6pt}

% ── دستورات کمکی ─────────────────────────────────────────
\newcommand{\placeholder}[1]{%
  \textcolor{mci-gold}{\textbf{[[ #1 ]]}}%
}

\newcommand{\speakercard}[4]{%
  % #1 نام  #2 سمت  #3 دانشگاه  #4 موضوع سخنرانی
  \begin{mdframed}[
    linecolor=mci-border,
    backgroundcolor=mci-light,
    roundcorner=4pt,
    innertopmargin=8pt,
    innerbottommargin=8pt,
    innerleftmargin=10pt,
    innerrightmargin=10pt,
    linewidth=1pt
  ]
  {\bfseries\color{mci-blue} #1}\\[2pt]
  {\small\color{mci-gray} #2 \textbar\ #3}\\[4pt]
  {\small\itshape موضوع: #4}
  \end{mdframed}
  \vspace{4pt}
}

% ─────────────────────────────────────────────────────────
\begin{document}
\settextfont{XB Niloofar}[Script=Arabic]

% ════════════════════════════════════════════════════════
%  صفحه کاور
% ════════════════════════════════════════════════════════
\begin{titlepage}

% پس‌زمینه گرادیان
\begin{tikzpicture}[remember picture, overlay]
  % نوار بالایی
  \fill[mci-blue]
    (current page.north west) rectangle
    ([yshift=-6cm]current page.north east);
  % نوار پایینی باریک
  \fill[mci-accent]
    ([yshift=1.2cm]current page.south west) rectangle
    (current page.south east);
  % خطوط تزئینی «مدار کوانتومی»
  \foreach \i in {1,...,6}{
    \draw[mci-accent, opacity=0.15, line width=0.6pt]
      ([xshift=\i*1.8cm, yshift=-6cm]current page.north west)
      -- ([xshift=\i*1.8cm]current page.south west);
  }
  % دایره‌های کوانتومی نمادین
  \draw[mci-gold, opacity=0.25, line width=1pt]
    ([xshift=3cm, yshift=-3cm]current page.north east) circle (2.5cm);
  \draw[mci-gold, opacity=0.15, line width=0.8pt]
    ([xshift=3cm, yshift=-3cm]current page.north east) circle (4cm);
  % لوگو placeholder
  \node[anchor=north east] at ([xshift=-1.5cm, yshift=-0.8cm]current page.north east){%
    \placeholder{لوگو همراه اول}
  };
\end{tikzpicture}

\vspace*{1cm}

% عنوان اصلی روی نوار آبی
{\color{white}\fontsize{22}{30}\selectfont\bfseries
  \begin{center}
    رویداد تخصصی
  \end{center}
}

{\color{mci-gold}\fontsize{28}{38}\selectfont\bfseries
  \begin{center}
    مخابرات کوانتومی\\[6pt]
    و رمزنگاری پساکوانتومی
  \end{center}
}

\vspace{1.2cm}

{\color{white}\large
  \begin{center}
    Quantum Communications \& Post-Quantum Cryptography
  \end{center}
}

\vspace{3cm}

% کادر اطلاعات رویداد
\begin{center}
\begin{tikzpicture}
  \node[
    draw=mci-border,
    fill=white,
    rounded corners=6pt,
    inner sep=16pt,
    text width=12cm,
    align=right
  ]{
    \begin{tabular}{r l}
      \textbf{\color{mci-blue}نام رویداد:}   & \placeholder{نام کامل رویداد} \\[6pt]
      \textbf{\color{mci-blue}تاریخ:}        & \placeholder{تاریخ برگزاری (شمسی/میلادی)} \\[6pt]
      \textbf{\color{mci-blue}مکان:}         & سالن همایش‌های ساختمان ستاره ونک \\[6pt]
      \textbf{\color{mci-blue}برگزارکننده:}  & تیم تحقیق و فناوری — همراه اول \\
    \end{tabular}
  };
\end{tikzpicture}
\end{center}

\vfill

{\small\color{mci-gray}\center{این سند محرمانه است و صرفاً جهت بررسی داخلی تهیه شده است.}}

\end{titlepage}

% ════════════════════════════════════════════════════════
%  بخش ۱: مقدمه و اهداف
% ════════════════════════════════════════════════════════
\section{مقدمه و اهداف رویداد}

تیم تحقیق و فناوری همراه اول در نظر دارد یک رویداد تخصصی یک‌روزه در حوزه
\textbf{مخابرات کوانتومی} و \textbf{رمزنگاری پساکوانتومی (PQC)} برگزار نماید.
این رویداد با هدف ایجاد بستری برای تبادل دانش، برندینگ واحد تحقیقات، و معرفی
دستاوردهای علمی تیم طراحی شده است.

\subsection{اهداف اصلی}

\begin{itemize}[rightmargin=0.5cm, itemsep=4pt]
  \item برندینگ و معرفی تیم تحقیق و فناوری همراه اول در سطح جامعه علمی-فناوری کشور
  \item برگزاری سخنرانی‌های تخصصی با حضور اساتید برجسته دانشگاهی و پژوهشگران
  \item رونمایی رسمی از کتاب ترجمه‌شده توسط تیم: \placeholder{عنوان کتاب}
  \item ایجاد شبکه ارتباطی بین متخصصان صنعت و دانشگاه در حوزه امنیت کوانتومی
  \item آشنایی مخاطبان با الگوریتم‌های مقاوم در برابر کوانتوم و استانداردهای NIST PQC
\end{itemize}

\subsection{محورهای تخصصی رویداد}

\begin{center}
\begin{tikzpicture}[
  topic/.style={
    draw=mci-blue, fill=mci-light,
    rounded corners=5pt, text width=5.2cm,
    align=center, minimum height=1.2cm,
    font=\small\bfseries\color{mci-blue}
  },
  node distance=0.5cm and 0.4cm
]
  \node[topic] (t1) {مخابرات کوانتومی\\{\footnotesize\color{mci-gray}QKD \& Quantum Networks}};
  \node[topic, right=of t1] (t2) {رمزنگاری پساکوانتومی\\{\footnotesize\color{mci-gray}NIST PQC Standards}};
  \node[topic, below=of t1] (t3) {امنیت شبکه پساکوانتومی\\{\footnotesize\color{mci-gray}Network Security}};
  \node[topic, right=of t3] (t4) {الگوریتم‌های مقاوم\\{\footnotesize\color{mci-gray}Lattice, Hash-based, ...}};
\end{tikzpicture}
\end{center}

% ════════════════════════════════════════════════════════
%  بخش ۲: سخنرانان مدعو
% ════════════════════════════════════════════════════════
\newpage
\section{سخنرانان مدعو}

\speakercard
  {دکتر اقلیدس}
  {\placeholder{سمت / رشته تخصصی}}
  {\placeholder{وابستگی دانشگاهی}}
  {\placeholder{موضوع سخنرانی}}

\speakercard
  {دکتر جواد صالحی}
  {استاد تمام — مهندسی برق}
  {دانشگاه صنعتی شریف}
  {\placeholder{موضوع سخنرانی}}

\speakercard
  {علی رضاخانی}
  {\placeholder{سمت / مرتبه علمی}}
  {دانشگاه صنعتی شریف}
  {\placeholder{موضوع سخنرانی}}

\speakercard
  {محمد رضایی}
  {\placeholder{سمت / مرتبه علمی}}
  {دانشگاه صنعتی شریف}
  {\placeholder{موضوع سخنرانی}}

\speakercard
  {امیر یوسفی}
  {\placeholder{سمت / مرتبه علمی}}
  {EPFL — سوئیس}
  {\placeholder{موضوع سخنرانی}}

\speakercard
  {محمد لامعی رشتی}
  {\placeholder{سمت / مرتبه علمی}}
  {پژوهشگاه دانش‌های بنیادی (IPM)}
  {\placeholder{موضوع سخنرانی}}

% ════════════════════════════════════════════════════════
%  بخش ۳: برنامه زمانی — سناریوی اول (با نهار)
% ════════════════════════════════════════════════════════
\newpage
\section{برنامه زمانی — سناریوی اول (با پذیرایی نهار)}

\begin{center}
{\small\color{mci-gray}زمان کل: ۸:۳۰ تا ۱۷:۳۰}
\end{center}

\vspace{0.4cm}

\begin{tikzpicture}[
  slot/.style={
    draw=mci-border, fill=white,
    rounded corners=3pt,
    minimum height=0.9cm,
    text width=11cm,
    inner xsep=8pt, align=right,
    font=\small
  },
  break/.style={
    draw=mci-accent!40, fill=mci-light,
    rounded corners=3pt,
    minimum height=0.9cm,
    text width=11cm,
    inner xsep=8pt, align=right,
    font=\small\color{mci-gray}
  },
  timebox/.style={
    text width=2.2cm, align=left,
    font=\small\bfseries\color{mci-blue}
  }
]

\def\ypos{0}
\def\ystep{-1.1}

% ردیف اول — پذیرش
\node[timebox] at (0,\ypos) {۸:۳۰ – ۹:۰۰};
\node[break] at (7,\ypos) {ثبت‌نام و پذیرش شرکت‌کنندگان};

\pgfmathsetmacro\ypos{\ypos+\ystep}
% افتتاحیه
\node[timebox] at (0,\ypos) {۹:۰۰ – ۹:۳۰};
\node[slot] at (7,\ypos) {\textbf{مراسم افتتاحیه} — معرفی تیم تحقیق و فناوری همراه اول};

\pgfmathsetmacro\ypos{\ypos+\ystep}
% سخنرانی ۱
\node[timebox] at (0,\ypos) {۹:۳۰ – ۱۰:۱۵};
\node[slot] at (7,\ypos) {\textbf{سخنرانی اول} — \placeholder{موضوع} | \placeholder{سخنران}};

\pgfmathsetmacro\ypos{\ypos+\ystep}
% سخنرانی ۲
\node[timebox] at (0,\ypos) {۱۰:۱۵ – ۱۱:۰۰};
\node[slot] at (7,\ypos) {\textbf{سخنرانی دوم} — \placeholder{موضوع} | \placeholder{سخنران}};

\pgfmathsetmacro\ypos{\ypos+\ystep}
% استراحت صبح
\node[timebox] at (0,\ypos) {۱۱:۰۰ – ۱۱:۲۰};
\node[break] at (7,\ypos) {پذیرایی — استراحت};

\pgfmathsetmacro\ypos{\ypos+\ystep}
% سخنرانی ۳
\node[timebox] at (0,\ypos) {۱۱:۲۰ – ۱۲:۰۵};
\node[slot] at (7,\ypos) {\textbf{سخنرانی سوم} — \placeholder{موضوع} | \placeholder{سخنران}};

\pgfmathsetmacro\ypos{\ypos+\ystep}
% نهار
\node[timebox] at (0,\ypos) {۱۲:۰۵ – ۱۳:۳۰};
\node[break] at (7,\ypos) {ناهار و نماز};

\pgfmathsetmacro\ypos{\ypos+\ystep}
% رونمایی کتاب
\node[timebox] at (0,\ypos) {۱۳:۳۰ – ۱۴:۱۵};
\node[slot,draw=mci-gold,fill=mci-gold!8] at (7,\ypos)
  {\textbf{\color{mci-gold}رونمایی از کتاب} — \placeholder{عنوان کتاب}};

\pgfmathsetmacro\ypos{\ypos+\ystep}
% سخنرانی ۴
\node[timebox] at (0,\ypos) {۱۴:۱۵ – ۱۵:۰۰};
\node[slot] at (7,\ypos) {\textbf{سخنرانی چهارم} — \placeholder{موضوع} | \placeholder{سخنران}};

\pgfmathsetmacro\ypos{\ypos+\ystep}
% سخنرانی ۵
\node[timebox] at (0,\ypos) {۱۵:۰۰ – ۱۵:۴۵};
\node[slot] at (7,\ypos) {\textbf{سخنرانی پنجم} — \placeholder{موضوع} | \placeholder{سخنران}};

\pgfmathsetmacro\ypos{\ypos+\ystep}
% استراحت بعدازظهر
\node[timebox] at (0,\ypos) {۱۵:۴۵ – ۱۶:۰۰};
\node[break] at (7,\ypos) {پذیرایی — استراحت};

\pgfmathsetmacro\ypos{\ypos+\ystep}
% سخنرانی ۶
\node[timebox] at (0,\ypos) {۱۶:۰۰ – ۱۶:۴۵};
\node[slot] at (7,\ypos) {\textbf{سخنرانی ششم} — \placeholder{موضوع} | \placeholder{سخنران}};

\pgfmathsetmacro\ypos{\ypos+\ystep}
% اختتامیه
\node[timebox] at (0,\ypos) {۱۶:۴۵ – ۱۷:۳۰};
\node[slot,fill=mci-blue!8,draw=mci-blue] at (7,\ypos)
  {\textbf{\color{mci-blue}اختتامیه} — جمع‌بندی، پرسش و پاسخ، عکس یادگاری};

\end{tikzpicture}

% ════════════════════════════════════════════════════════
%  بخش ۴: برنامه زمانی — سناریوی دوم (بدون نهار)
% ════════════════════════════════════════════════════════
\newpage
\section{برنامه زمانی — سناریوی دوم (بدون پذیرایی نهار)}

\begin{center}
{\small\color{mci-gray}زمان کل: ۸:۳۰ تا ۱۳:۳۰}
\end{center}

\vspace{0.4cm}

\begin{tikzpicture}[
  slot/.style={
    draw=mci-border, fill=white,
    rounded corners=3pt, minimum height=0.9cm,
    text width=11cm, inner xsep=8pt, align=right,
    font=\small
  },
  break/.style={
    draw=mci-accent!40, fill=mci-light,
    rounded corners=3pt, minimum height=0.9cm,
    text width=11cm, inner xsep=8pt, align=right,
    font=\small\color{mci-gray}
  },
  timebox/.style={
    text width=2.2cm, align=left,
    font=\small\bfseries\color{mci-blue}
  }
]

\def\ypos{0}
\def\ystep{-1.1}

\node[timebox] at (0,\ypos) {۸:۳۰ – ۹:۰۰};
\node[break] at (7,\ypos) {ثبت‌نام و پذیرش شرکت‌کنندگان};

\pgfmathsetmacro\ypos{\ypos+\ystep}
\node[timebox] at (0,\ypos) {۹:۰۰ – ۹:۲۰};
\node[slot] at (7,\ypos) {\textbf{مراسم افتتاحیه} — معرفی تیم تحقیق و فناوری};

\pgfmathsetmacro\ypos{\ypos+\ystep}
\node[timebox] at (0,\ypos) {۹:۲۰ – ۱۰:۰۰};
\node[slot] at (7,\ypos) {\textbf{سخنرانی اول} — \placeholder{موضوع} | \placeholder{سخنران}};

\pgfmathsetmacro\ypos{\ypos+\ystep}
\node[timebox] at (0,\ypos) {۱۰:۰۰ – ۱۰:۴۰};
\node[slot] at (7,\ypos) {\textbf{سخنرانی دوم} — \placeholder{موضوع} | \placeholder{سخنران}};

\pgfmathsetmacro\ypos{\ypos+\ystep}
\node[timebox] at (0,\ypos) {۱۰:۴۰ – ۱۱:۰۰};
\node[break] at (7,\ypos) {پذیرایی — استراحت};

\pgfmathsetmacro\ypos{\ypos+\ystep}
\node[timebox] at (0,\ypos) {۱۱:۰۰ – ۱۱:۴۰};
\node[slot] at (7,\ypos) {\textbf{سخنرانی سوم} — \placeholder{موضوع} | \placeholder{سخنران}};

\pgfmathsetmacro\ypos{\ypos+\ystep}
\node[timebox] at (0,\ypos) {۱۱:۴۰ – ۱۲:۲۰};
\node[slot] at (7,\ypos) {\textbf{سخنرانی چهارم} — \placeholder{موضوع} | \placeholder{سخنران}};

\pgfmathsetmacro\ypos{\ypos+\ystep}
\node[timebox] at (0,\ypos) {۱۲:۲۰ – ۱۳:۰۰};
\node[slot,draw=mci-gold,fill=mci-gold!8] at (7,\ypos)
  {\textbf{\color{mci-gold}رونمایی از کتاب} — \placeholder{عنوان کتاب}};

\pgfmathsetmacro\ypos{\ypos+\ystep}
\node[timebox] at (0,\ypos) {۱۳:۰۰ – ۱۳:۳۰};
\node[slot,fill=mci-blue!8,draw=mci-blue] at (7,\ypos)
  {\textbf{\color{mci-blue}اختتامیه} — جمع‌بندی، پرسش و پاسخ};

\end{tikzpicture}

% ════════════════════════════════════════════════════════
%  بخش ۵: معرفی کتاب
% ════════════════════════════════════════════════════════
\newpage
\section{معرفی کتاب}

\begin{mdframed}[
  linecolor=mci-gold,
  linewidth=2pt,
  backgroundcolor=mci-gold!5,
  roundcorner=6pt,
  innertopmargin=14pt,
  innerbottommargin=14pt,
  innerleftmargin=14pt,
  innerrightmargin=14pt
]

\begin{minipage}{0.2\textwidth}
  \centering
  \placeholder{تصویر جلد}
\end{minipage}
\hfill
\begin{minipage}{0.76\textwidth}
  \textbf{\large\color{mci-blue}\placeholder{عنوان فارسی کتاب}}\\[4pt]
  {\color{mci-gray}\placeholder{عنوان انگلیسی / اصلی کتاب}}\\[8pt]
  \textbf{مترجمان:} \placeholder{نام مترجمان}\\[4pt]
  \textbf{ناشر:} \placeholder{نام ناشر}\\[4pt]
  \textbf{سال انتشار:} \placeholder{سال}\\[8pt]
  {\small این کتاب توسط اعضای تیم تحقیق و فناوری همراه اول ترجمه شده
  و یکی از منابع اصلی آموزشی در حوزه \placeholder{موضوع کتاب} محسوب می‌شود.}
\end{minipage}

\end{mdframed}

% ════════════════════════════════════════════════════════
%  بخش ۶: اطلاعات مخاطبان هدف و مشخصات اجرایی
% ════════════════════════════════════════════════════════
\section{مخاطبان هدف}

\begin{itemize}[rightmargin=0.5cm, itemsep=3pt]
  \item متخصصان حوزه مخابرات و امنیت اطلاعات در شرکت‌های فناوری
  \item دانشجویان و پژوهشگران دانشگاهی در رشته‌های مرتبط
  \item مدیران و تصمیم‌گیرندگان حوزه امنیت سایبری
  \item علاقه‌مندان به فناوری‌های نوظهور کوانتومی
\end{itemize}

\section{مشخصات اجرایی}

\begin{tabularx}{\textwidth}{>{\bfseries\color{mci-blue}}r X}
  \toprule
  مکان برگزاری & سالن همایش‌های ساختمان ستاره ونک \\
  \midrule
  ظرفیت تخمینی & \placeholder{تعداد نفر} نفر \\
  \midrule
  زبان رویداد & فارسی (با اسلایدهای انگلیسی) \\
  \midrule
  ضبط و پخش & \placeholder{آنلاین / فقط حضوری} \\
  \midrule
  هزینه ثبت‌نام & \placeholder{رایگان / مبلغ} \\
  \bottomrule
\end{tabularx}

% ════════════════════════════════════════════════════════
%  صفحه پایانی
% ════════════════════════════════════════════════════════
\newpage
\thispagestyle{empty}

\begin{tikzpicture}[remember picture, overlay]
  \fill[mci-blue]
    (current page.north west) rectangle
    ([yshift=-4cm]current page.north east);
  \fill[mci-accent]
    ([yshift=1.2cm]current page.south west) rectangle
    (current page.south east);
\end{tikzpicture}

\vspace*{0.8cm}
{\color{white}\Large\bfseries\center{با سپاس از توجه شما}}

\vspace{2cm}

\begin{center}
\begin{mdframed}[
  linecolor=mci-border,
  backgroundcolor=mci-light,
  roundcorner=6pt,
  innertopmargin=16pt, innerbottommargin=16pt,
  innerleftmargin=20pt, innerrightmargin=20pt
]
  {\large\bfseries\color{mci-blue} تیم تحقیق و فناوری}\\[4pt]
  {\color{mci-gray} شرکت ارتباطات سیار ایران (همراه اول)}\\[12pt]
  \textbf{ایمیل:} \placeholder{آدرس ایمیل تماس}\\[4pt]
  \textbf{وبسایت:} \placeholder{آدرس وبسایت}\\[4pt]
  \placeholder{لوگو همراه اول}
\end{mdframed}
\end{center}

\end{document}

